\section{Related Work}
\label{sec:related-work}

Blockchain-based consent management for healthcare data sharing has been studied extensively, with smart contracts used to record, enforce, and audit consent across institutions~\cite{kakarlapudi2021systematic, welzel2025enabling,li2024you}.
These works focus on making consent decisions verifiable and auditable, and the privacy mechanisms surveyed in this domain mainly protect health data contents and user identities~\cite{li2025privacy}.
The confidentiality of the governing consent policy, and what public-chain enforcement artifacts reveal about it, is typically out of scope.

Beyond consent management, related work addresses individual requirements of policy-private consent enforcement (policy confidentiality, metadata-leakage control, verifiable auditability, or a stable verification interface), but none of the approaches reviewed here combines them on a public chain.
Encrypted data-sharing mechanisms, such as CP-ABE and PRE~\cite{bethencourt2007ciphertext,ateniese2006improved}, control access to encrypted data. However, as noted in Section~\ref{subsec:crypto-access}, standard CP-ABE exposes the policy itself, while PRE-based delegation may reveal authorization patterns that data encryption is not designed to hide.
Transaction-privacy systems, including Zerocash-style protocols~\cite{sasson2014zerocash} and Ethereum-oriented systems such as ZETH and AZTEC~\cite{rondelet2019zeth,williamson2018aztec}, provide public verification over private transaction witnesses, but they mainly target asset ownership, values, or participant relationships rather than compliance with a provider-defined consent policy.

Hidden-policy CP-ABE schemes reduce direct policy exposure by hiding the access structure in the ciphertext~\cite{yang2024hamed,ding2026badsabe}.
Wu et al.~\cite{wu2023blockchain} instead realize attribute-based access control (ABAC) by encrypting policy and attribute values under an additively homomorphic cryptosystem.
In this scheme, a smart contract evaluates access conditions over encrypted values with assistance from each organization's node, which provides partial decryption results and zero-knowledge proofs of their correctness.
However, the Boolean structure of the policy and its attribute names remain visible on the ledger, and per-attribute match results can be recovered from the ciphertexts and partial decryptions recorded with each decision~\cite{wu2023blockchain}.
The scheme also operates on a consortium blockchain, where enforcement artifacts are not globally visible as they are on a public blockchain.
Other approaches provide broader forms of confidentiality.
Functional encryption reveals function outputs while hiding inputs~\cite{boneh2011functional}, but does not by itself provide publicly verifiable authorization evidence.
Private-state blockchain systems, including Secret Network and Oasis/Ekiden, protect smart-contract state or execution through encrypted state and trusted execution environments (TEEs)~\cite{secretnetwork_docs,cheng2019ekiden}.
Policy or state confidentiality alone, however, does not ensure that public-chain enforcement artifacts reveal no additional policy information.
TPS-ZKPI therefore focuses on limiting policy information exposed during public-chain consent enforcement while preserving public verifiability of each authorization decision.

Self-sovereign identity (SSI) architectures give holders control over their verifiable claims, including the ability to determine which verifiers may access them.
Buccafurri et al.~\cite{buccafurri2024holder} support this control through ciphertext-policy attribute-based proxy re-encryption. The holder stores encrypted claims off-chain under a self-defined attribute policy, and a semi-trusted, non-colluding consensus provider re-encrypts them for verifiers whose certified attributes satisfy that policy. A smart contract notarizes hashes of storage, policy-change, and access operations so that they can later be attributed to their parties, and a trusted key-generation authority issues all attribute keys.
Zero-knowledge credential presentations likewise let a patient authenticate anonymously on public or private chains~\cite{madine2025zero}, and surveys of zero-knowledge proofs for blockchain privacy catalogue this design space~\cite{shashidhara2024promise,zhou2024leveraging}.
TPS-ZKPI focuses on protecting the provider's consent policy during authorization.
The requester's claims serve as private witness inputs to a reusable zero-knowledge circuit that verifies satisfaction of the committed policy.

Privacy-preserving access-control and policy-evaluation mechanisms have been studied in credential-based settings built on the eXtensible Access Control Markup Language (XACML) and the Security Assertion Markup Language (SAML)~\cite{ardagna2009xacml}, as well as in policy evaluation by multi-party computation (MPC) using homomorphic encryption (HE) and secure function evaluation~\cite{sheikhalishahi2019privacy}.
These approaches support confidential authorization computation, but they are not primarily designed for settings where policy-evaluation artifacts (calldata, public inputs, emitted events, timing, and gas usage) are permanently and globally visible on a public blockchain.
In addition, these approaches do not directly provide compact evidence that a smart contract can verify against a committed, provider-defined policy instance.

Reusable proving systems, including universal-setup succinct non-interactive arguments of knowledge (SNARKs), composable proof systems, and zero-knowledge virtual machines (zkVMs)~\cite{gabizon2019plonk,campanelli2019legosnark,liu2025ceno}, improve maintainability by reducing the need for policy-specific circuits or verifier contracts.
However, reusable proving alone does not guarantee policy privacy during public-chain consent enforcement.
Policy-dependent program representations, public inputs, calldata layouts, logs, or gas profiles may still reveal information about the hidden policy.
This motivates TPS-ZKPI's use of a domain-specific reusable policy interpreter, which preserves proof reuse while exposing only a policy-independent public verification interface.

Studies of metadata leakage on public chains show that cryptographic privacy
guarantees may be weakened by observable transaction-level artifacts such as
timing and volume patterns~\cite{beres2021blockchain}. Techniques such as batching,
decoys, and traffic shaping can reduce some observable signals~\cite{miller2023vtbc}, but they are
usually designed for payment privacy, anonymity sets, or transaction-linkability
settings. They do not directly provide hidden-policy compliance verification
or auditable consent enforcement. In contrast, TPS-ZKPI treats metadata leakage as
part of the policy-privacy objective rather than as a separate transaction-level concern.

Table~\ref{tab:related-work-gap} summarizes the reviewed approaches by policy confidentiality, metadata-leakage mitigation, and verifiable auditability, and notes the scope and limitations of each.
Partial indicates limited support or a need for application-specific design;
N/A denotes a criterion that does not apply to the cited approach.


\begin{table}[htbp]
	\centering
	\caption{Positioning of representative approaches against policy-private consent
		enforcement under public-chain observability.}
	\label{tab:related-work-gap} \footnotesize
	% Column widths sum to <1 so that the inter-column glue fits inside the text block.
	\begin{tabular}{@{}>{\raggedright\arraybackslash}p{0.19\linewidth}>{\raggedright\arraybackslash}p{0.09\linewidth}>{\raggedright\arraybackslash}p{0.09\linewidth}>{\raggedright\arraybackslash}p{0.09\linewidth}p{0.38\linewidth}@{}}
		\hline
		\textbf{Approach}                                                                                                                                     & \textbf{Policy confidentiality} & \textbf{Metadata-leakage mitigation} & \textbf{Verifiable auditability} & \textbf{Scope and limitations}                                                                                                                                                                \\
		\hline
		Blockchain consent management~\cite{kakarlapudi2021systematic,welzel2025enabling}                                                                     & No                              & N/A                                  & Yes                              & Records and enforces consent verifiably on chain, but policy confidentiality and policy-dependent leakage are typically out of scope.                                                         \\
		Data / transaction privacy~\cite{bethencourt2007ciphertext,ateniese2006improved,hu2023access,sasson2014zerocash,rondelet2019zeth,williamson2018aztec} & No                              & Partial                              & Partial                          & Protects encrypted data access or transaction semantics, but not the governing consent policy or policy-dependent authorization workflow.                                                     \\
		Hidden-policy ABE~\cite{yang2024hamed,ding2026badsabe}                                                                                                & Partial                         & Partial                              & Partial                          & Reduces direct policy exposure, but may leak through workflows, updates, or policy-holder interactions.                                                                                       \\
		Homomorphic hidden-policy ABAC~\cite{wu2023blockchain}                                                                                                & Partial                         & N/A                                  & Yes                              & Hides policy and attribute values on a permissioned ledger, but the Boolean structure, attribute names, and per-attribute match results remain recoverable from it.                           \\
		SSI verifier authorization~\cite{buccafurri2024holder}                                                                                                & No                              & Partial                              & Partial                          & Releases the holder's claims only to attribute-qualified verifiers and notarizes access for later attribution, but protects the requester's claims rather than the provider's consent policy. \\
		MPC / HE policy evaluation~\cite{sheikhalishahi2019privacy}                                                                                           & Yes                             & N/A                                  & N/A                              & Supports confidential policy evaluation off-chain; the cited design does not specify public-chain enforcement or authorization auditing.                                                      \\
		Universal SNARKs / zkVMs~\cite{gabizon2019plonk,campanelli2019legosnark,liu2025ceno}                                                                  & Partial                         & Partial                              & Yes                              & Supports reusable proving and public verification; policy hiding and a fixed verification interface require application-specific design.                                                      \\
		Metadata-leakage defenses~\cite{miller2023vtbc}                                                                                                       & N/A                             & Yes                                  & Partial                          & Mitigates transaction timing and volume leakage and supports verifiable execution, but does not directly provide hidden-policy consent enforcement.                                           \\
		TPS-ZKPI                                                                                                                                              & Yes (under TPS trust)           & Yes (modeled channels)               & Yes                              & Supports bounded Boolean policies; confidentiality requires trust in the TPS.                                                                                                                 \\
		\hline
	\end{tabular}
\end{table}
